\clearpage
\thispagestyle{empty}
\pdfbookmark[0]{Beobachtungen der magnetischen Variation in Göttingen und Leipzig. Ausgabe und Quellen}{oct1834-edition}
\begin{center}\Large\textbf{Beobachtungen der magnetischen Variation\\in Göttingen und Leipzig}\end{center}

Die vollständige Fassung der Werkausgabe von Carl Friedrich Gauß’ Bericht über die Beobachtungen am 1. und 2. Oktober 1834 ist eine KI-Übertragung aus den \textit{Werken}, Band V (Göttingen, 1867), S.~525--528, in modernes Deutsch. Der Satz wurde mit KI erstellt. Der Text enthält das Rechenbeispiel mit den sechs Ablesungen, den fünf Zwischenmitteln und dem abschließenden Mittelwert.

Die Zeitschriftenfassung erschien als Beitrag XLII in Johann Christian Poggendorffs \textit{Annalen der Physik und Chemie}, Band~33 (109), 1834, S.~426--433. Zu dieser Fassung gehören die beiden größeren Beobachtungstabellen, der Tafelverweis, die Herausgeberanmerkungen und der ausführlichere Terminplan. Maßgeblich für die Übertragung sind Text und Umfang der Werkausgabe.

\medskip\noindent
\href{https://gdz.sub.uni-goettingen.de/id/PPN236006339}{Deutsche Werkausgabe (1867)}\par
\noindent\href{https://www.digitale-sammlungen.de/view/bsb10130346?page=440}{Originaler Zeitschriftenbeitrag (1834)}\par
\noindent\href{https://doi.org/10.1007/978-3-642-49319-5_29}{Deutscher Nachdruck (1877)}
\clearpage
